%LTeX: language=es-AR

\documentclass[12pt]{amsart}

%Fuente
\usepackage[utf8]{inputenc}
\usepackage[spanish]{babel}
\usepackage[T1]{fontenc}
\usepackage{lmodern}
\usepackage{epigraph}
\usepackage{lipsum}
\usepackage[a4paper, margin=2.5cm]{geometry}
\usepackage{multicol}
\usepackage{amsmath}
\usepackage{amssymb}
\usepackage[shortlabels]{enumitem}
\usepackage{listings}
\lstdefinestyle{micode}{
    basicstyle=\ttfamily\small,
    breaklines=true,
    columns=fullflexible,
    keepspaces=true,
    showstringspaces=false
}

\DeclareMathOperator{\mcd}{mcd}

\renewcommand{\refname}{Bibliografía}

%diagramas
\usepackage{tikz-cd}

%enumerar páginas y encabezados
\pagestyle{headings}

%Comandos para escribir C más rápido
\usepackage{mathrsfs}
\newcommand{\CC}{\mathbb{C}}
\newcommand{\NN}{\mathbb{N}}
\newcommand{\QQ}{\mathbb{Q}}
\newcommand{\RR}{\mathbb{R}}
\newcommand{\ZZ}{\mathbb{Z}}
\newcommand{\FF}{\mathbb{F}}
\newcommand{\cont}{\mathcal{C}}

\newcommand{\pp}{\mathfrak{p}}

\newcommand{\aaa}{\mathfrak{a}}
\newcommand{\bbb}{\mathfrak{b}}

\newcommand{\OO}{\mathcal{O}}

\newcommand{\leg}[2]{\left( \frac{#1}{#2} \right)}
\newcommand{\minimat}[4]{\left(\begin{smallmatrix} #1 & #2 \\ #3 & #4 
\end{smallmatrix}\right)}
\newcommand{\lp}{\left(}
\newcommand{\rp}{\right)}
\newcommand{\lb}{\left|}
\newcommand{\rb}{\right|}
\newcommand{\lc}{\left<}
\newcommand{\rc}{\right>}
\newcommand{\lco}{\left[}
\newcommand{\rco}{\right]}

%dobles implicaciones
\usepackage{csquotes}

%Entornos con estilo 'tipo teorema'
\theoremstyle{plain}
\newtheorem{teo}[equation]{Teorema}
\newtheorem{ej}[equation]{Ejercicio}
\newtheorem{prop}[equation]{Proposición}
\newtheorem{lem}[equation]{Lema}
\newtheorem{coro}[equation]{Corolario}
\newtheorem{defi}[equation]{Definición}
\newtheorem{obs}[equation]{Observación}
\newtheorem*{nota*}{Notación}
\renewcommand*{\proofname}{Demostración}

% Símbolo de Tate--Shafarevich
\DeclareFontFamily{U}{wncy}{}
\DeclareFontShape{U}{wncy}{m}{n}{<->wncyr10}{}
\DeclareSymbolFont{cyrletters}{U}{wncy}{m}{n}
\DeclareMathSymbol{\Sha}{\mathord}{cyrletters}{"58}

%Hipervínculos cliqueables, en color azul
\usepackage[colorlinks=true,allcolors=blue]{hyperref}

\title{Título}

\author{Pedro Nicolás Peña}

\begin{document}

\maketitle

% con Alt + Z te acomoda las lineas

\section*{Sección}

Vamos a trabajar mucho con la curva $E_{t^2}: y^2 = x^3 + t^2$ porque tiene un punto de 3 torsión, permitiéndonos 3-descenso. Tiene grupo de torsión isomorfo a $\ZZ/3\ZZ$ sobre $\QQ$ (porque no te queda otra) \textbf{ver sobre $\QQ(\omega)$}. Si $t = \frac{r}{s}$ entonces queda isomorfa a $E_{rs}: y^2 = x^3 +r^2s^4$. Las curvas  de la forma $y^2 = x^3 + d$ se llaman curvas de Mordell y hay mucha bibliografia sobre ellas.

Para un ejemplo de rango 3 ver https://arxiv.org/pdf/1604.02693.

%For curves with rational 3-torsion, your student could instead use this paper by Feng et al. (https://www.math.clemson.edu/~kevja/PAPERS/Selmer_Groups_3-Torsion-2011-11-10.pdf), which provides formulas for the 3-Selmer group (that should simplify considerably once you restrict to j=0). Alternatively, Chan’s recent work on the 3-isogeny Selmer groups of y^2 = x^3 + n^2 could be helpful. 

% Note that bounding the rank typically requires computing both the phi-Selmer group and the phi'-Selmer group (where phi is the 3-isogeny and phi' is its dual). However, for E_{k^2} over Q(sqrt(-3)), the curve is isomorphic to its dual E_{-27k^2}. I have not worked out the details, but I think this should allow your student to bound the rank using only the phi-Selmer group

\textbf{La  curva isógena:}

La isogenia con kernel generado por $(0,t)$ nos da la curva $E':y^2 = x^3 -27t^2$

\textbf{Los puntos de torsión:}

El punto $P=(0,t)$ es de orden $3$, y es el único sobre $\QQ$. Sobre $\QQ(\omega)$ puede haber otro si el parametro $t$ es cuatro veces un cubo.




% Sí ponemos $P = ()$ y $Q = (r(u-r), iru(u-r))$ con $u^2=r^2-s^2$

% \[
% \begin{array}{|c|c|c|c|c|}
% \hline
% +   & \OO    & P      & 2*P    & 3*P \\ \hline

% \OO & \OO     & (rs,rs(r+s)) & (0,0)     & (rs,-rs(r+s)) \\ \hline

% Q& (-r(r-u),  & (-s(s-iu), & (-r(r+u), & (-s(s+iu), \\

% &  -iru(r-u)) & su(s-iu))  & iru(r+u)) & -su(s+iu)) \\ \hline

% 2*Q& (-r^2,0) & (-rs,-rs(r-s))& (-s^2,0)& (-rs,rs(r-s)) \\ \hline

% 3*Q & (-r(r-u),& (-s(s+iu, & (-r(r+u), & (-s(s-iu), \\ 

% &   iru(r-u)) & su(s+iu)) & -iru(r+u)) & -su(s-iu)) \\ \hline
% \end{array}
% \]

% (esta tabla está chequeada a mano con Sage)

\textbf{Discriminante:}

\[
\Delta_{E_{t^2}} = -432t^4 = -2^4 3^3 t^4
\]

Su factorización en $\QQ(\omega)$ es: $2^4 \omega^6 t^4$ si $t\equiv 2 \pmod 3$ primo de $\QQ$ y $2^4 \omega^6 t_1^4 t_2^4$ si $t\equiv 1 \pmod 3$ primo de $\QQ$.

\textbf{j-invariante:}

0

\textbf{Notas del paper de Ben:}

Since an elliptic curve $E_\alpha/\QQ(i)$ of the form (1.2) has complex multiplication by $\ZZ[i]$, the group $E_\alpha(\QQ(i))$ is a $\ZZ[i]$-module. In particular, the rank of $E_\alpha$ is necessarily even. Thus,
to prove $E_\alpha$ has rank 2, it suffices to produce a single non-torsion $\QQ(i)$-point and show that
the rank is at most 2. To this end, we first compute all of the $\QQ(i)$-torsion points. We then
use descent via a 2-isogeny to bound the rank from above.

\textbf{Notas para seguir el texto de Cohen y Pazuki:}

Antes que nada no olvidar que las curvas que estamos viendo son isomorfas a su dual sobre $\QQ(\sqrt{-3})$

Su $b$ es nuestro $t^2$ y $b=b_1b_3^3$ con $b_1$ libre de cubos

$\alpha$ va a ser nuestro mapa que toma puntos de la curva y nos da elementos de $K^*/{(K^*)}^3$ con $\alpha(\OO)=1$, $\alpha((0, t)) = 1/(2t^2)$ y $\alpha((x, y)) = y-t^2$.

$\hat\alpha$ va a ser lo mismo desde la isógena, y nos queda $\hat\alpha(\OO)=1$, $\hat\alpha((0, -t)) = 1/(54t^2)$ y $\hat\alpha((x, y)) = y+27t^2$.

La joda teórica es $Ker(\alpha) = im(\hat\phi)$. También $|Im(\alpha)||Im(\hat\alpha)| = 3^{r+1}$. El corolario es 
\[
\frac{E(\QQ)}{\hat\phi(E'(\QQ))} \simeq Im(\alpha)
\]

Para $\QQ$ vale lo siguiente: $\bar{u}\in \QQ^\times / {(\QQ^\times)}^3$ esta en la imagen de $\alpha$ si y solo si tiene un repr en $u\in \QQ$ tal que la siguiente tiene solución no trivial en $\ZZ$ (o $\QQ$, es lo mismo):
\[
uX^3 + (1/u)Y^3 + 2bZ^3 -2aXYZ= 0
\]

Y eso se puede simplificar: para una solución tomá $u=u_1^2u_2$ (troll) todo libre de cuadrados y $u_1\perp u_2$, y entonces $u_1u_2 | 2b$ y la solubilidad es equiv a la de:
\[
u_1X^3 + u_2Y^3 + (2b/(u_1u_2))Z^3 - 2aXY Z = 0 
\]

Osea para las que tenemos (esto es en $\QQ$ no olvidar):

\[
u_1X^3 + u_2Y^3 + (2t^2/(u_1u_2))Z^3 = 0
\]

\[
u_1X^3 + u_2Y^3 - (54t^2/(u_1u_2))Z^3 = 0
\]

Obvio para los casos de $\OO$ y los puntos de torsión tenemos soluciones y son:
\[
\alpha(\OO) = 1 = u \implies (X,Y,Z) = (1,-1,0)
\]
\[
\alpha((0, t)) = 1/2t = u \implies (X,Y,Z) = (0,1,-1)
\]

Y el caso $y-t^2 = uz^3$ tiene solución $(X,Y,Z) = (z^2,-x,z)$

Dada una solución podemos buscar el punto del que venía.



\textbf{Selmer:}

Ver que onda $Sel(E(K))$ para $K\neq \QQ$

\textbf{Como ver ${\QQ(\sqrt{-3})}^\times / {({\QQ(\sqrt{-3})}^\times)}^3$:}


Note that since we are working over $Q(\sqrt{-3})$, $p = 3$ is the only ramified prime. If $p \equiv 2 \pmod 3$, then $p$ is an inert prime.
And if $p \equiv 1 \pmod 3$, then $p$ is a split prime



\textbf{Puntos libres:}

https://personal.math.ubc.ca/~bennett/BeGa-data.html


\begin{prop}
Para una curva elíptica definida sobre $K$, una extensión cuadrática $L=K(\sqrt{D})$ y el twist $E_D$ vale

\[
rg(E(L)) = rg(E(K)) + rg(E_D(K))
\]
\end{prop}

\textbf{A parametrised family of Mordell curves:}

Arranca con $v_i^2 - k^2= u_i^3$ para $i=1,2,3$ y con eso se va a armar 3 soluciones independientes (con $x=u$ y $v=y$)

Si todos los $u_i$ tienen un factor común $m$ lo elimina y te quedan dos ecuaciones: pongamos $a=u_1/m$, $b=u_2/m$, $c=u_3/m$

\[
b^3(v_1^2-k^2) = a^3 (v_2^2-k^2)
\]
\[
c^3(v_1^2-k^2) = a^3 (v_3^2-k^2)
\]

Despues pone $v_i = w_it+k$ e igualando en $t$ tenes una cubica en $w_i$ de la que sacas varias soluciones.



\textbf{Que se sabe:}

Hay infinitas curvas de rango $r$ para $0\leq r \leq 4$ sobre todo cuerpo de números. Las da Zywina en \cite{ZywinaSmallRank} que todas tienen full 2-torsion (generada por $(0,0)$)(obvio podes matarla y tenes de torsion trivial) con $\Sha [2] = 0$ (falta leer un poco mas de ese). Todo vino de una secuencia de papers: \cite{ZywinaRankNotJump}, \cite{ZywinaRankTwo} y \cite{ZywinaRankOneStability}. Usa numeros de Tamagawa para calcular los simbolos de Kodaira y el global root number porque (la conjetura de paridad dice que) te da la paridad del rango, y estas cosas te dan restricciones sobre las especializaciones.

De Elkies hay curvas con rango alto, todas sobre $\QQ$, y ordenadas por grupo de torsion: Para el trivial el record es $29$ y está en un post random, y para varios grupos ($\ZZ /n\ZZ$ para $n\in \{2,3,4,6,7\}$) hay un paper: \cite{ElkiesKlagsbrun2020}. Lo hace con superficies K3.















\end{document}